\section{Geometric meaning of solutions}

Linear systems have an algebraic form, but they also have a geometric meaning. This geometric interpretation is often the easiest way to understand why a system may have one solution, no solution, or infinitely many solutions.

The central idea is simple: each equation describes a set of points. The solution set of the system is the intersection of all these sets.

\subsection*{Two variables: lines in the plane}

In two variables, a linear equation
\[
ax+by=c
\]
usually describes a line in the plane. A system of two linear equations describes two lines, and a solution is a point that lies on both lines.

\begin{examplebox}
Consider
\[
\begin{cases}
x+y=5,\\
2x-y=1.
\end{cases}
\]
The first equation describes one line. The second equation describes another line. The pair \((2,3)\) is a solution because the point \((2,3)\) lies on both lines.
\end{examplebox}

This is the geometric meaning of solving the system: find the common point of the two lines.

\subsection*{One solution}

If two lines in the plane intersect in one point, the system has exactly one solution.

\begin{examplebox}
The system
\[
\begin{cases}
x+y=5,\\
2x-y=1
\end{cases}
\]
has exactly one solution. Solving gives
\[
x=2,
\qquad
y=3.
\]
Geometrically, the two lines meet at the point \((2,3)\).
\end{examplebox}

When a system has exactly one solution, the conditions are independent enough to determine the unknowns completely.

\subsection*{No solution}

If two lines are parallel and distinct, they never meet. Then the system has no solution.

\begin{examplebox}
Consider
\[
\begin{cases}
x+y=1,\\
x+y=4.
\end{cases}
\]
The left-hand sides are the same, but the right-hand sides are different. The same expression \(x+y\) cannot equal both \(1\) and \(4\) at the same time.

Geometrically, these are two parallel distinct lines. Therefore the system has no solution.
\end{examplebox}

A system with no solution is called inconsistent. The equations impose conditions that cannot all be satisfied simultaneously.

\subsection*{Infinitely many solutions}

If two equations describe the same line, then every point on that line is a solution. The system has infinitely many solutions.

\begin{examplebox}
Consider
\[
\begin{cases}
x+y=3,\\
2x+2y=6.
\end{cases}
\]
The second equation is just twice the first equation. It gives the same condition, not a new independent condition.

Therefore the solution set is the whole line
\[
x+y=3.
\]
There are infinitely many solutions, for example \((0,3)\), \((1,2)\), and \((3,0)\).
\end{examplebox}

This case is not a failure of mathematics. It means that the equations do not contain enough independent information to determine one point.

\subsection*{Three variables: planes in space}

In three variables, a linear equation
\[
ax+by+cz=d
\]
usually describes a plane in space. A system of linear equations describes the common points of several planes.

Several things may happen:
\begin{itemize}
\item three planes may meet at one point,
\item they may have no common point,
\item they may share a line,
\item they may all be the same plane.
\end{itemize}

The algebraic classification remains the same: one solution, no solution, or infinitely many solutions. The geometry becomes richer because planes can intersect in more ways than lines in the plane.

\subsection*{Reading geometry from algebra}

The algebraic form of a system often reveals the geometric situation. For example, proportional left-hand sides may indicate parallel or identical lines. A contradiction such as
\[
0=5
\]
means that the conditions cannot all be satisfied. A row of zeros such as
\[
0=0
\]
means that one equation became redundant.

These patterns will appear naturally during Gaussian elimination.

\begin{summarybox}
When interpreting a system geometrically, ask:
\begin{itemize}
\item What geometric object does each equation describe?
\item Is the solution set an intersection of lines, planes, or other linear objects?
\item Do the objects meet in one point?
\item Are they parallel or inconsistent?
\item Is one condition a repetition of another condition?
\item Does the geometry suggest one solution, no solution, or infinitely many solutions?
\end{itemize}
\end{summarybox}

The geometric viewpoint does not replace algebraic methods. It explains what those methods are detecting.